\documentclass{gescons}

\genre {Entrevista}
\author{Cesar Machado}
\title{Grafoproéxis: A~Escrita Conscienciológica enquanto Megafoco Existêncial}
% \paginaurl{https://www.youtube.com/live/oJfBRY6j8w4}

\begin{document}
    \makeentrevistatitle
    \coverart{back/Cesar_Machado}
    \vspace*{\fill}

    \begin{multicols}{2}

\begin{center}
    \includegraphics[width=6cm]{articles/entrevista/mockups/Cesar_Machado}
\end{center}


\textbf{1. Qual foi a~motivação para a~escrita da obra? Por que a~definição deste tema para publicação de um livro?}

Eu já apresentava uma relação com a~escrita desde a~pré-adolescência, quando tinha 13 anos de idade e,~na disciplina de Educação Artística, em que o~objetivo era fazer algum tipo de desenho, escrevi um pequeno conto infantil (comento sobre isso no livro). Posteriormente, acabei me dedicando à~Medicina enquanto profissão. Mas a~escrita realmente surgiu na condição de necessidade íntima, a~partir do contato com as técnicas de escrita conscienciológica, desenvolvidas e~aplicadas pelo professor Waldo Vieira. Foi aí que percebi a~importância da escrita na minha proéxis. Associada à~pesquisa de retropersonalidade no campo da autoria de livros, surgiu a~ideia do tema Grafoproéxis, tendo a~escrita conscienciológica enquanto megafoco existencial.

\textbf{2. Quais foram as principais percepções, intra e~extrafísicas, durante a~pesquisa e~a~escrita da obra? E~posterior ao lançamento?}

Foram muitas ocorrências relacionadas à~escrita do livro, também relato isso na obra. Principalmente sincronicidades com datas, nomes, números, livros, peças teatrais, filmes, etnias, dentre outras. Sincronicidades retrocognitivas e~também relacionadas ao conteúdo do livro. Também ocorreram sincronicidades após o~livro entrar em fase final de revisão, o~que reforçou a~importância do tema. Ainda ocorreram fenômenos de clarividência com ex-autores em dinâmica parapsíquica, além de percepções dentro do processo de parapsiquismo impressivo em dinâmicas e~no dia a~dia da pesquisa e~escrita da obra.


\textbf{3. Qual o~maior aprendizado com a~escrita desta obra?}

Houve muito contrafluxo, pressão extrafísica e~ocorrências intrafísicas que dificultaram o~término da obra, após o~primeiro parecer na Editares, em 2019. Para mim, então, o~maior aprendizado foi a~importância da resiliência e~da autodefesa energética na escrita da obra conscienciológica. Neste livro, especificamente, fiz a~evocação de muitas personalidades do passado, em épocas diferentes, e~muitas coisas em contextos que acabaram não sendo positivos na história da humanidade, tendo por pano de fundo os escritores e~seus escritos. Este livro tem esse objetivo, de realizar acertos grupocármicos no campo autoral e~desensinar o~que foi ensinado equivocadamente em vidas humanas anteriores.

% \begin{pullquote}
% ``O maior aprendizado foi a~importância da resiliência e~da autodefesa energética na escrita da obra conscienciológica.''
% \end{pullquote}

\textbf{4. O~que poderia dizer como incentivo para que mais pesquisadores invistam na publicação de obras conscienciológicas?}

Vale lembrar que a~Conscienciologia é~uma ciência e,~desse modo, a~escrita é~parte essencial para o~seu crescimento e~desenvolvimento. Além disso, a~escrita conscienciológica é~uma excelente forma de realizar acertos grupocármicos e,~também, funcionar enquanto cápsula do tempo, visando o~autorrevezamento para as próximas existências intrafísicas. Considero estes os principais incentivos para os pesquisadores que desejam investir na escrita de livros conscienciológicos. Importa enfatizar também 2 objetivos centrais desta obra, considerando o~papel do escritor para a~maxiproéxis conscienciológica: agir ao modo de agente retrocognitor (ativando as memórias dos Cursos Intermissivos) e~atrator ressomático (atraindo novos escritores \emph{intermissivistas-grafoproexistas}).

    
    \end{multicols}
    \vspace*{\fill}
\end{document}
